\documentclass[a4paper,12pt]{article}

\usepackage[utf8]{inputenc}
\usepackage[portuguese]{babel}
\usepackage{fancyhdr}
\usepackage{amsmath}
\usepackage{amsthm}
\usepackage{amssymb}
\usepackage[portuguese,algoruled,lined,linesnumbered]{algorithm2e}

% ---- PREENCHA ----
\newcommand{\Lista}{2}
\newcommand{\Nome}{Leonardo Heidi Almeida Murakami}
\newcommand{\NUSP}{11260186}
% -----------------------------

\pagestyle{fancy}
\fancyhf{}
\fancyhead[L]{\Nome}
\fancyhead[R]{NUSP \NUSP}
\fancyfoot[L]{Lista \Lista}
\fancyfoot[C]{MAC0338}

\begin{document}

\section*{Exercício 1}
\subsection*{Parte a}
\subsubsection*{Probabilidade}
As comparações da linha 6 serão executadas sempre que v[i] <= max{v[1]..v[i-1]}, e, como a distribuição é uniforme conforme o enunciado sabemos que a probabilidade de, na iteração i de v[i] ser maior que v[1..i-1] será de $\frac{1}{i}$

A probabilidade entao de executarmos a linha 6 será o complemento desta $P(\text{linha 6}) = 1 - \frac{1}{i}$
\subsubsection*{Esperanca}
Temos portanto que, a quantidade esperada de execucoes/comparacoes da linha 6 será
\begin{align}
    E[X_6]
    &= \sum_{i=2}^n 1 - \sum_{i=2}^n \frac{1}{i} \\
    &= n - H_n
\end{align}

\subsection*{Parte b}
A probabilidade de a linha 7 ser executada é exatamente a probabilidade de v[i] ser o valor minimo de v[1..i], que, na distribuicao uniforme, é de $\frac{1}{i}$, portanto
\begin{align}
    E[X_7]
    &= \sum_{i=2}^n \frac{1}{i} \\
    &= (\frac{1}{2} + \frac{1}{3} + \dots + \frac{1}{n}) \\
    &= (1 + \frac{1}{2} + \frac{1}{3} + \dots + \frac{1}{n}) - 1 \\
    &= H_n - 1
\end{align}
\section*{Exercício 2}
Vamos testar a proposta para o caso base:
\begin{align}
    M(0)
    &\geq \frac{1}{2}(0 + 1)\lg{(0 + 1)} \\
    &\geq \frac{1}{2}\lg{1} \\
    &\geq 0
\end{align}
Que é verdadeiro, para M(1)
\begin{align}
    M(0)
    &\geq \frac{1}{2}(1 + 1)\lg{(1 + 1)} \\
    &\geq \frac{2}{2}\lg{2} \\
    &\geq 1
\end{align}
Que é verdadeiro também.

\section*{Exercício 6}
O quicksort nao utiliza arrays auxiliares durante sua execucao, portanto todo o seu consumo de espaco vem da parte de dividir e conquistar. Disso sabemos que, no pior caso, caso o algoritmo sempre escolha o pior pivo possivel ele tera que fazer n chamadas recursivas ate que enfim a lista esteja particionada ordenadamente. Logo, O(n)

\section*{Exercício 7}
\section{Exercício 9}
\begin{algorithm}[H]
\SetKwProg{Def}{def}{:}{}
\SetAlgoLined
\caption{median}

\KwData{duas listas ordenadas $X[1..N]$ e $Y[1..n]$}
\KwResult{a mediana das duas listas se elas fossem intercaladas}
\Def{median(X, Y, p, q, r, s)}{
    $medianX = p + floor((q - p)/2)$\;
    $medianY = r + floor((s - r)/2)$\;
    \If{$p == q$}{
        return $min(X[medianX],Y[medianY])$\;
    }
    \If{$X[medianX] \le Y[medianY]$}{
        return median(X, Y, medianX, q, r, medianY)\;
    }
    \Else{
        return median(X, Y, p, medianX, medianY, s)\;
    }
}
\end{algorithm}
\end{document}